\section{Background and Related Work}
\label{sec:related}

\subsection{PTM-Isoform Drug Design (PTMI-DD)}
\label{sec:related-ptmidd}

The PTMI-DD framework targets the disease-relevant PTM-isoform $\poiptm$ rather than the wild-type protein $\poiwt$. Selectivity is gained because $\poiptm$ carries chemistry and conformation that $\poiwt$ does not. Concrete instantiations cluster into four operational variants: (a)~covalent ligation at non-conserved residues exposed by the PTM (ibrutinib at the BTK cysteine), (b)~allosteric binding in a PTM-induced or PTM-stabilized pocket (asciminib at the myristoyl pocket of Bcr-Abl), (c)~modulation of PTM-dependent protein-protein interactions (phospho-tyrosine recognition by SH2 domains; methyl-K9 read by HP1$\gamma$), and (d)~PTM-conditioned proteasomal degradation via E3 recruitment (PROTACs anchored on CRBN/VHL/MDM2/IAP). The framework synthesizes a strong demand-side story but does not solve the prospective nomination question: given an arbitrary candidate POI, which PTM-isoform and which E3-PROTAC handle will yield a tractable degrader? Our work attacks the fourth (PROTAC) variant of this demand. The other three remain open for future work in the same operational style.

\subsection{Ternary-Complex Scoring for PROTAC Degraders}
\label{sec:related-ternary}

Three machine-learning ternary-complex predictors define the current state of the art on the standard CRBN/VHL benchmarks. DeepTernary~\cite{deepternary-2025} introduces a graph-neural-network scorer trained on TernaryDB, predicting the pTernary score for each $\tuple$ tuple. PROTAC-STAN~\cite{protac-stan-2025} uses a structure-aware transformer with explicit recruitment-geometry priors. ET-PROTACs~\cite{et-protacs-2024} couples an equivariant network with a PROTAC-DB-derived training set. All three score canonical UniProt sequences for the POI: a PTM-isoform $\poiptm$ and its wild-type counterpart $\poiwt$ are scored identically up to per-tuple noise. The training data itself is PTM-blind: TernaryDB and PROTAC-DB do not annotate POI PTM state at the time of recruitment. Our work treats these scorers as black boxes and adds two layers \emph{outside} them: PTM-conditioned POI structure prediction (\S\ref{sec:method-pipeline}) and per-POI noise-floor calibration (\S\ref{sec:method-calibration}). The pipeline is scorer-agnostic, so improvements in the underlying scorer flow through.

\subsection{PTM-Conditioned Structure Prediction}
\label{sec:related-structure}

AlphaFold~3~\cite{af3-2024} and Boltz-2 ingest native CCD-PTM tokens for phospho-Ser/Thr/Tyr, mono/di/tri-methyl-Lys, ubiquitin attachment, and several other modifications. Both are single-state predictors: a single inference returns one structure, not a conformational ensemble, even though many PTMs drive disorder-to-order transitions or population shifts. The AF3 paper itself flags this limit (cereblon predicted in the closed state when given as apo); recent work~\cite{bouvier-2025} finds that AF3-phospho only modestly improves over AF2 for phospho-driven conformational changes. We take the conservative position that single-state PTM-occupied prediction is sufficient for ranking only when its uncertainty is absorbed by an external calibration. The per-POI noise-floor protocol (\S\ref{sec:method-calibration}) is that external calibration, and the MD-relaxed fallback route (\S\ref{sec:method-pipeline}) supplies an independent structural prior when CCD-PTM token coverage is missing. The complementary line of work that attempts to predict full PTM-conditioned conformational ensembles~\cite{af3-2024}---a harder problem---is orthogonal to ours: when ensemble prediction matures, this paper's pipeline can consume the ensemble mode of $\poiptm$ without architectural change.

\subsection{E3-Substrate Predictors and the Recruitable-E3 Axis}
\label{sec:related-e3}

Existing E3-substrate prediction tools, of which UbiBrowser~\cite{ubibrowser-2017} is the most cited, score candidate $(\eligase, \mathrm{substrate})$ pairs from sequence, domain, and network features. These predictors answer ``which E3 ubiquitylates this substrate?'' rather than ``which E3-PROTAC handle recruits a PTM-isoform?''---they operate on physiological substrates, not on bifunctional-degrader recruitment. The recruitable-E3 axis for PROTAC design is empirically narrow: $\approx 80\%$ of PROTAC-DB entries use CRBN or VHL, with MDM2 and IAP family accounting for most of the remainder. Extending the recruitable axis requires PTM-aware POI-side ranking, which is the gap our method closes for phospho-degrons (\S\ref{sec:experiments-main}). The methylation-track failure (\S\ref{sec:experiments-scope}) shows the axis is not extended uniformly across PTM types.

\subsection{Position of This Work}

We position $\dpt$ as a method-paper contribution: it is the simplest scorer-agnostic, noise-floor-calibrated PTM-aware nomination axis that we know of. The point of departure from prior art is that we do not modify the underlying ternary scorer (no retraining, no fine-tuning) and we do not require a working PTM-conditioned ensemble predictor (single-state $\poiptm$ structures plus a noise-floor bound suffice). The pre-registered multi-track benchmark, including the deliberate negative result on methylation, is the contribution on the evaluation side.
